% AA-52 — skończone podgrupy SL2(C).
\chapter{Binarne grupy wielościenne i skończone podgrupy SL(2)}
\chaptermark{Binarne grupy wielościenne}
\label{aa:ch:binary-polyhedral}

W rozdziale~\ref{aa:ch:finite-pgl2} poznaliśmy typy
skończonych podgrup grupy rzutowej.
Teraz odtworzymy grupy macierzowe leżące nad nimi.
Istotna różnica polega na tym, że nie każdy obrót rzędu
dwa ma podniesienie rzędu dwa: w grupie specjalnej
liniowej jego podniesienia mają rząd cztery.

\section{Sprowadzenie do grupy unitarnej}

\begin{twierdzenie}[Unitarność po zmianie bazy]
\label{aa:tw:binary-unitary}
Dla każdej skończonej podgrupy
$H\le\operatorname{SL}_2(\mathbb C)$ istnieje
$S\in\operatorname{SL}_2(\mathbb C)$ takie, że
$SHS^{-1}\subset\operatorname{SU}(2)$.
\end{twierdzenie}
\begin{proof}
Niech $A=\sum_{h\in H}h^*h$, gdzie $h^*$ oznacza
macierz sprzężoną hermitowsko. Macierz $A$ jest dodatnio
określona, gdyż dla $v\ne0$ zachodzi
$v^*Av=\sum_{h\in H}\|hv\|^2>0$.
Dla $g\in H$ mamy
$g^*Ag=\sum_{h\in H}(hg)^*(hg)=A$.
Istnieje dodatnio określony hermitowski pierwiastek
$P=A^{1/2}$, otrzymany przez diagonalizację unitarną
$A$ i zastąpienie jej dodatnich wartości własnych
ich dodatnimi pierwiastkami.
Wówczas $(PgP^{-1})^*(PgP^{-1})=I$.
Wyznacznik $PgP^{-1}$ wynosi $1$.
Aby uzyskać macierz zmiany bazy o wyznaczniku $1$,
wystarczy przyjąć $S=(\det P)^{-1/2}P$,
co nie zmienia sprzężenia.
\end{proof}

Z Lematu~\ref{aa:lem:pgl2-rotations} mamy surjekcję
$\pi:\operatorname{SU}(2)\to\operatorname{SO}(3)$,
daną sprzężeniem na rzeczywistej przestrzeni
hermitowskich macierzy o śladzie zero.
Jej jądro to $\{I,-I\}$.
Przez $\rho:\operatorname{SL}_2(\mathbb C)
\to\operatorname{PGL}_2(\mathbb C)$ oznaczamy
homomorfizm rzutowy. Jego ograniczenie do
$\operatorname{SU}(2)$ utożsamia się z $\pi$
za pomocą izomorfizmu z poprzedniego rozdziału.

\begin{lemat}[Elementy rzędu dwa i pełne przeciwobrazy]
\label{aa:lem:binary-full-lift}
Niech $H\le\operatorname{SL}_2(\mathbb C)$ będzie skończona.
Jedynym elementem rzędu $2$ w $\operatorname{SL}_2(\mathbb C)$
jest $-I$. Jeżeli $|H|$ jest parzyste, to
$H=\rho^{-1}(\rho(H))$. Jeśli $|H|$ jest nieparzyste,
to $H$ jest cykliczna.
\end{lemat}
\begin{proof}
Jeśli $M^2=I$, to $M$ jest diagonalizowalna, ponieważ
$T^2-1$ ma różne pierwiastki. Jej wartości własne
należą do $\{1,-1\}$. Warunek $\det M=1$ wymusza
ich równość, więc $M=\pm I$.
Dla parzystego $|H|$ twierdzenie Cauchy'ego daje $-I\in H$.
Jeśli $\rho(M)\in\rho(H)$, istnieje $h\in H$ takie, że
$Mh^{-1}\in\ker\rho=\{\pm I\}\subset H$,
więc $M\in H$.
Dla nieparzystego $|H|$ ograniczenie $\rho|_H$ jest
iniektywne. Z klasyfikacji w
Twierdzeniu~\ref{aa:tw:pgl2-classification} wynika,
że jedyne grupy rzutowe o nieparzystym rzędzie są cykliczne.
\end{proof}

\section{Grupy cykliczne i dwucykliczne}

\begin{definicja}[Grupa dwucykliczna]
\label{aa:def:binary-dicyclic}
Dla całkowitego $n\ge2$ połóżmy
$\zeta=\exp(\pi i/n)$ oraz
\[
 A=\begin{pmatrix}\zeta&0\\0&\zeta^{-1}\end{pmatrix},
 \qquad J=\begin{pmatrix}0&1\\-1&0\end{pmatrix}.
\]
Grupę $\operatorname{Dic}_n=\langle A,J\rangle
\le\operatorname{SU}(2)$ nazywamy grupą dwucykliczną
(binarna grupa dwuścienna).
\end{definicja}

\begin{twierdzenie}[Postać i relacje grupy dwucyklicznej]
\label{aa:tw:binary-dicyclic}
Dla każdego $n\ge2$ grupa $\operatorname{Dic}_n$
ma rząd $4n$, obraz rzutowy izomorficzny z $D_n$
rzędu $2n$ oraz prezentację
\[
 \langle a,j\mid
 a^{2n}=1,\ j^2=a^n,\ jaj^{-1}=a^{-1}\rangle.
\]
\end{twierdzenie}
\begin{proof}
Bezpośrednie mnożenie macierzy daje
$A^{2n}=I$, $J^2=-I=A^n$ oraz $JAJ^{-1}=A^{-1}$.
Relacje pozwalają zapisać każdy wyraz w generatorach
i ich odwrotnościach w postaci $a^k j^\varepsilon$,
gdzie $0\le k<2n$ i $\varepsilon\in\{0,1\}$:
przesuwamy $j$ na prawo przez $ja^k=a^{-k}j$,
a każdą parę $j^2$ zastępujemy przez $a^n$.
Stąd grupa zadana prezentacją ma najwyżej $4n$ elementów.
Macierze $A^k$ są parami różne i diagonalne,
zaś $A^kJ$ są parami różne i mają zera na przekątnej.
Otrzymujemy $4n$ różnych macierzy, zatem naturalny
surjektywny homomorfizm z grupy zadanej prezentacją
do grupy macierzy jest izomorfizmem.
Iloraz przez $\{\pm I\}$ jest generowany przez klasy
$A,J$ o rzędach $n,2$ z relacją odwracania; ma $2n$
elementów, więc jest grupą dwuścienną $D_n$.
\end{proof}

\begin{przyklad}[Grupa kwaternionowa]
\label{aa:prz:binary-q8}
Dla $n=2$ mamy $A=\operatorname{diag}(i,-i)$.
Macierze $A,J,AJ$ mają kwadrat $-I$, a ponadto
$AJ=-JA$. Osiem elementów $\pm I,\pm A,\pm J,\pm AJ$
tworzy grupę kwaternionową $Q_8$.
Iloraz przez $\{\pm I\}$ jest grupą Kleina.
To pokazuje, że grupa nad $D_2$ nie jest grupą
$C_2\times D_2$: ta ostatnia ma siedem inwolucji,
podczas gdy $Q_8$ ma jedną.
\end{przyklad}

\begin{lemat}[Przeciwobrazy grup cyklicznych i dwuściennych]
\label{aa:lem:binary-cyclic-dihedral}
Niech $G\le\operatorname{PGL}_2(\mathbb C)$ będzie
skończona. Jeżeli $G\simeq C_n$, to jej pełny
przeciwobraz względem $\rho$ jest sprzężony
z grupą diagonalną rzędu $2n$.
Jeżeli $G\simeq D_n$ dla $n\ge2$, to jej pełny
przeciwobraz jest sprzężony z $\operatorname{Dic}_n$.
Sprzężenia można wybrać w $\operatorname{SL}_2(\mathbb C)$.
\end{lemat}
\begin{proof}
Dla $n=1$ pierwsza grupa to $\{\pm I\}$.
Dla $n>1$ wybieramy dwa punkty stałe generatora
i przenosimy je do $0,\infty$. Grupa rzutowa ma
wówczas postać $z\mapsto\eta^k z$, $\eta$ pierwotny
pierwiastek $n$-tego stopnia z jedności. Jej pełny
przeciwobraz to macierze
$\operatorname{diag}(\lambda,\lambda^{-1})$ z
$\lambda^{2n}=1$.
W przypadku dwuściennym istnieje orbita dwupunktowa
na mocy danych stabilizatorów z rozdziału 51.
Po przeniesieniu jej do $0,\infty$ grupa rzutowa jest
generowana przez $z\mapsto\eta z$ i $z\mapsto c/z$,
$c\ne0$. Sprzężenie przez $z\mapsto dz$ zastępuje
$c$ przez $d^2c$; wybieramy $d^2=c^{-1}$.
Pełny przeciwobraz jest teraz generowany przez
$A,J$ z Definicji~\ref{aa:def:binary-dicyclic}, gdyż
ich obraz rzutowy generuje $G$, a $-I=A^n$.
Każde użyte sprzężenie rzutowe ma reprezentanta
o wyznaczniku $1$ po pomnożeniu przez odpowiedni skalar.
\end{proof}

\section{Binarne grupy wielościenne}

\begin{definicja}[Binarne grupy obrotów]
\label{aa:def:binary-exceptional}
Niech $T,O,I_{\mathrm{rot}}\le\operatorname{SO}(3)$
oznaczają grupy obrotów czworościanu, sześcianu
i dwudziestościanu z przykładów
w rozdziale~\ref{aa:ch:finite-pgl2}.
Definiujemy
$2T=\pi^{-1}(T)$, $2O=\pi^{-1}(O)$ oraz
$2I=\pi^{-1}(I_{\mathrm{rot}})$.
Nazywamy je odpowiednio binarną grupą czworościenną,
ośmiościenną i dwudziestościenną.
\end{definicja}

Surjektywność $\pi$ i jej jądro dają dokładne ciągi
\[
 \begin{aligned}
 1&\to C_2\to 2T\to A_4\to1,\\
 1&\to C_2\to 2O\to S_4\to1,\\
 1&\to C_2\to 2I\to A_5\to1.
 \end{aligned}
\]
Grupy mają odpowiednio $24,48,120$ elementów.
Żaden z tych ciągów nie rozszczepia się:
sekcja homomorficzna wysyłałaby inwolucję ilorazu
na inwolucję różną od $-I$, ponieważ $-I$ należy
do jądra. Sprzeczność z Lematem~\ref{aa:lem:binary-full-lift}.
Ten sam argument dotyczy ciągu dla $\operatorname{Dic}_n$.

\begin{lemat}[Rzędy podniesień]
\label{aa:lem:binary-orders}
Niech $g\in\operatorname{SO}(3)$ ma rząd $m\ge1$.
Jego dwa przeciwobrazy $u,-u\in\operatorname{SU}(2)$
mają dla parzystego $m$ oba rząd $2m$,
a dla nieparzystego $m$ jeden rząd $m$, drugi $2m$.
\end{lemat}
\begin{proof}
Dla $m=1$ podniesienia to $\pm I$.
Dla $m>1$ po zmianie osi obrót ma kąt $2\pi k/m$,
gdzie $\gcd(k,m)=1$. Na mocy konstrukcji surjekcji
$\pi$ podniesienia są sprzężone z
$\pm\operatorname{diag}(e^{-\pi i k/m},e^{\pi i k/m})$.
Dla każdego podniesienia jego rząd jest wielokrotnością
$m$, a jego $m$-ta potęga należy do $\{\pm I\}$,
więc rząd wynosi $m$ lub $2m$.
Jeśli $m$ jest parzyste, $k$ jest nieparzyste i obie
$m$-te potęgi wynoszą $-I$. Jeśli $m$ jest nieparzyste,
zmiana znaku podniesienia zmienia znak $m$-tej potęgi,
co daje dokładnie jeden element każdego z dwóch rzędów.
\end{proof}

\begin{przyklad}[Rozkłady rzędów elementów]
\label{aa:prz:binary-order-distributions}
Z rozkładów permutacji na rozłączne cykle otrzymujemy
liczby elementów rzędów $1,2,3,4,5$ w grupach
$A_4,S_4,A_5$:
\[
\begin{array}{c|rrrrr}
 &1&2&3&4&5\\ \hline
 A_4&1&3&8&0&0\\
 S_4&1&9&8&6&0\\
 A_5&1&15&20&0&24
\end{array}
\]
Na przykład inwolucji w $A_5$ jest $5\cdot3=15$
(wybór punktu stałego i parowania czterech pozostałych).
Cykli długości $3$ jest $\binom53\cdot2=20$,
a cykli długości $5$ jest $4!=24$.
Z Lematu~\ref{aa:lem:binary-orders} wynika tabela
dla grup binarnych:
\[
\begin{array}{c|rrrrrrr}
 &1&2&3&4&5&6&8\\ \hline
 2T&1&1&8&6&0&8&0\\
 2O&1&1&8&18&0&8&12\\
 2I&1&1&20&30&24&20&0
\end{array}
\]
W ostatnim wierszu należy dodatkowo uwzględnić
$24$ elementy rzędu $10$. Sumy wynoszą $24,48,120$.
\end{przyklad}

\section{Jednoznaczność z dokładnością do sprzężenia}

\begin{lemat}[Charakter reprezentacji obrotów]
\label{aa:lem:binary-rotation-character}
Niech $K$ będzie jedną z grup $A_4,S_4,A_5$.
Jeżeli $r_1,r_2:K\to\operatorname{SO}(3)$ są iniektywnymi
homomorfizmami, to po ewentualnym złożeniu $r_2$
z automorfizmem $K$ ich zespolone reprezentacje
trójwymiarowe mają ten sam charakter.
\end{lemat}
\begin{proof}
Macierz obrotu ma wartości własne $1,e^{it},e^{-it}$.
Dla elementu rzędu $m$ kąt ma postać $2\pi k/m$,
gdzie $\gcd(k,m)=1$. Ślad wynosi $1+2\cos(2\pi k/m)$.
Dla $m=1,2,3,4$ ślad zależy wyłącznie od $m$:
odpowiednio $3,-1,0,1$. To wystarcza dla $A_4,S_4$.

W $A_5$ cykle długości pięć tworzą dwie klasy
sprzężoności po $12$ elementów. Istotnie centralizator
cyklu $c$ długości pięć w $S_5$ jest $\langle c\rangle$:
permutacja komutująca z $c$ jest określona przez obraz
jednego punktu. Wszystkie jego elementy są parzyste,
więc centralizator w $A_5$ ma rząd $5$.
Każda klasa ma zatem $60/5=12$ elementów.
Elementy $c,c^2$ nie są sprzężone w $A_5$:
po oznaczeniu punktów przez $\mathbb Z/5\mathbb Z$
i przyjęciu $c:x\mapsto x+1$ każde sprzężenie $c$ z $c^2$
ma postać $x\mapsto 2x+b$. Mnożenie przez $2$
jest czterocyklem na elementach niezerowych,
a przesunięcie jest parzyste, więc taki sprzęgający
element jest nieparzysty.
Ślady $r_\ell(c)$ i $r_\ell(c^2)$ są różne:
wynoszą w pewnej kolejności
$1+2\cos(2\pi/5)$ oraz $1+2\cos(4\pi/5)$.
Sprzężenie $K=A_5$ przez nieparzystą permutację
daje automorfizm zamieniający te dwie klasy
(gdyby zachowywał klasę $c$, istniałby nieparzysty
element centralizatora $c$).
Możemy więc uzgodnić wartości obu charakterów
również na elementach rzędu pięć.
\end{proof}

\begin{lemat}[Od równoważności zespolonej do sprzężenia obrotów]
\label{aa:lem:binary-real-intertwiner}
Niech $K$ będzie skończoną grupą,
a $r_1,r_2:K\to\operatorname{SO}(3)$ homomorfizmami.
Jeśli ich zespolone reprezentacje mają ten sam charakter,
to istnieje $Q\in\operatorname{SO}(3)$ spełniające
$Qr_1(k)Q^{-1}=r_2(k)$ dla każdego $k\in K$.
\end{lemat}
\begin{proof}
Twierdzenie o charakterach reprezentacji skończonych
nad $\mathbb C$ (rozdział 42) daje odwracalną
zespoloną macierz $Z$ taką, że
$Zr_1(k)=r_2(k)Z$.
Zapiszmy $Z=X+iY$ dla macierzy rzeczywistych $X,Y$.
Obie spełniają te same równania liniowe, ponieważ
$r_1(k),r_2(k)$ są rzeczywiste.
Wielomian $p(t)=\det(X+tY)$ nie jest zerowy,
gdyż $p(i)=\det Z\ne0$. Istnieje zatem $t\in\mathbb R$
takie, że $L=X+tY$ jest odwracalna.
Niech $D=L^{\mathsf T}L$.
Z równania $Lr_1(k)=r_2(k)L$ i ortogonalności
otrzymujemy $r_1(k)^{\mathsf T}Dr_1(k)=D$,
a zatem $D$ komutuje z każdym $r_1(k)$.
Każda przestrzeń własna dodatnio określonej
macierzy symetrycznej $D$ jest przez nie zachowywana;
więc również $D^{-1/2}$ komutuje z $r_1(k)$.
Macierz $Q_0=LD^{-1/2}$ jest ortogonalna
i nadal spełnia równania przeplatania.
Jeśli jej wyznacznik jest ujemny, zastępujemy ją
przez $-Q_0$, co w wymiarze trzy zmienia znak wyznacznika,
ale nie zmienia sprzężenia.
\end{proof}

\begin{twierdzenie}[Klasyfikacja skończonych podgrup specjalnej grupy liniowej]
\label{aa:tw:binary-classification}
Każda skończona podgrupa $\operatorname{SL}_2(\mathbb C)$
jest sprzężona w tej grupie z jedną z następujących:
\[
\begin{gathered}
 C_m^{\mathrm{diag}}=
 \{\operatorname{diag}(\lambda,\lambda^{-1}):\lambda^m=1\},
 \quad m\ge1,\\
 \operatorname{Dic}_n\quad(n\ge2),\qquad 2T,\quad 2O,\quad 2I.
\end{gathered}
\]
Typy w tej liście są parami różne z dokładnością
do izomorfizmu (z wyjątkiem powtórzenia identycznego parametru).
\end{twierdzenie}
\begin{proof}
Na mocy Twierdzenia~\ref{aa:tw:binary-unitary}
możemy przyjąć $H\subset\operatorname{SU}(2)$.
Jeżeli $|H|$ jest nieparzyste, grupa jest cykliczna
na mocy Lematu~\ref{aa:lem:binary-full-lift}; jej generator
jest diagonalizowalny unitarnie, a macierz sprzęgającą
można pomnożyć przez skalar, by miała wyznacznik $1$.
Otrzymujemy $C_{|H|}^{\mathrm{diag}}$.
Jeśli $|H|$ jest parzyste, jest pełnym przeciwobrazem
swego obrazu rzutowego. Dla obrazu cyklicznego lub
dwuściennego stosujemy Lemat~\ref{aa:lem:binary-cyclic-dihedral}.
W pozostałych przypadkach obraz $\pi(H)$ jest
izomorficzny z jedną z grup $K=A_4,S_4,A_5$.
Wybieramy izomorfizmy $K\to\pi(H)$ i $K\to G_0$,
gdzie $G_0$ to odpowiednia grupa obrotów
z Definicji~\ref{aa:def:binary-exceptional}.
Lematy~\ref{aa:lem:binary-rotation-character}
i~\ref{aa:lem:binary-real-intertwiner} dają
$Q\in\operatorname{SO}(3)$ sprzęgające te obrazy
(jako podgrupy; automorfizm $K$ nie zmienia obrazu).
Wybierzmy $U\in\pi^{-1}(Q)$. Wówczas
$UHU^{-1}=\pi^{-1}(G_0)$, ponieważ obie strony
są pełnymi przeciwobrazami $G_0$.
To dowodzi klasyfikacji z dokładnością do sprzężenia.

Grupy cykliczne odróżniamy od pozostałych przez
przemienność i cykliczność. Grupy dwucykliczne
są nieprzemienne dla $n\ge2$, a ich rzędy wyznaczają $n$.
Grupy $2T,2O,2I$ mają różne rzędy.
Gdy grupa dwucykliczna ma jeden z tych rzędów,
ma element rzędu odpowiednio $12,24,60$ (generator
podgrupy diagonalnej). Tabela z
Przykładu~\ref{aa:prz:binary-order-distributions}
wyklucza te rzędy dla odpowiednich grup wielościennych.
\end{proof}

\section{Naturalna reprezentacja i równanie Riccatiego}

\begin{wniosek}[Nieprzywiedlność i prymitywność]
\label{aa:wn:binary-primitive}
Niech $H\le\operatorname{SL}_2(\mathbb C)$ będzie skończona.
Jej naturalna reprezentacja na $\mathbb C^2$ jest
przywiedlna dokładnie dla grup cyklicznych,
nieprzywiedlna imprymitywna dokładnie dla grup
$\operatorname{Dic}_n$ ($n\ge2$) i prymitywna dokładnie
dla $2T,2O,2I$ (z dokładnością do sprzężenia).
\end{wniosek}
\begin{proof}
Grupa cykliczna ma wspólną prostą własną.
Jeżeli $H$ ma prostą niezmienniczą, po sprowadzeniu
do grupy unitarnej jej ortogonalne dopełnienie jest
również niezmiennicze. W odpowiedniej bazie $H$ jest
podgrupą macierzy diagonalnych wyznacznika $1$,
więc zanurza się w $\mathbb C^\times$ i jest cykliczna.
Dla grup dwucyklicznych macierze $A,J$ zachowują
nieuporządkowaną parę osi współrzędnych, lecz grupa
nie jest cykliczna. Pozostałe trzy grupy są
nieprzywiedlne; zachowanie pary prostych dawałoby
orbitę rzutową o co najwyżej dwóch elementach.
Najmniejsze długości orbit, obliczone w
rozdziale~\ref{aa:ch:finite-pgl2}, wynoszą odpowiednio
$4,6,12$. Zatem ich reprezentacje są prymitywne.
\end{proof}

\begin{stwierdzenie}[Wielomian związany z orbitą prostych]
\label{aa:st:binary-orbit-form}
Niech $H\le\operatorname{SL}_2(\mathbb C)$ będzie skończona,
a $\mathcal O$ orbitą działania $H$ na
$\mathbb P^1(\mathbb C)$. Dla każdej prostej
$L\in\mathcal O$ wybierzmy niezerowy liniowy funkcjonał
$\ell_L$ na $\mathbb C^2$ o jądrze $L$.
Wielomian jednorodny
$F_{\mathcal O}(v)=\prod_{L\in\mathcal O}\ell_L(v)$
ma stopień $|\mathcal O|$. Istnieje charakter
jednowymiarowy $\chi:H\to\mathbb C^\times$ taki, że
$F_{\mathcal O}(h^{-1}v)=\chi(h)F_{\mathcal O}(v)$
dla $h\in H$ i $v\in\mathbb C^2$.
Jeśli $-I\in H$, to $\chi(-I)=(-1)^{|\mathcal O|}$.
\end{stwierdzenie}
\begin{proof}
Funkcjonał $\ell_L\circ h^{-1}$ ma jądro $hL$,
więc jest niezerową wielokrotnością $\ell_{hL}$.
Ponieważ $h$ permutuje orbitę, iloczyn tych
funkcjonałów jest skalarną wielokrotnością
$F_{\mathcal O}$. Skalar jest jednoznaczny,
bo wielomian nie jest zerowy.
Działanie $(h\cdot F)(v)=F(h^{-1}v)$ jest reprezentacją
na przestrzeni wielomianów jednorodnych;
porównanie działań $h_1h_2$ oraz $h_1,h_2$ pokazuje
$\chi(h_1h_2)=\chi(h_1)\chi(h_2)$.
Ostatnia równość wynika z jednorodności wielomianu.
\end{proof}

\begin{przyklad}[Para osi dla grupy dwucyklicznej]
Dla $\operatorname{Dic}_n$ orbita osi współrzędnych
daje wielomian $F(x,y)=xy$.
Działanie $A$ zachowuje $F$, zaś
$J^{-1}(x,y)=(-y,x)$ daje $(J\cdot F)(x,y)=-xy$.
Wielomian nie jest więc niezmiennikiem, lecz rozpina
jednowymiarową podprzestrzeń niezmienniczą.
Jego kwadrat jest już niezmiennikiem.
\end{przyklad}

Dla grup $2T,2O,2I$ konstrukcja orbit daje
wielomiany półniezmiennicze stopni $4,6,12$.
Zauważmy, że istnienie takiej podprzestrzeni
w potędze symetrycznej reprezentacji sprzężonej
nie jest równoznaczne z niezmienniczością każdego
wektora tej podprzestrzeni.
W przypadku skończonej prymitywnej różniczkowej
grupy Galois równania $z''=qz$ z $q\in\mathbb C(x)$
Twierdzenie~\ref{aa:tw:binary-classification} pozwala
zidentyfikować pełną grupę macierzową, nie tylko jej
obraz rzutowy. Stopnie orbit są zgodne z parametrami
trzeciego przypadku algorytmu Kovacica.
Nie wyprowadzamy z samej klasyfikacji pełnych
warunków lokalnych algorytmu ani tablic charakterów
grup binarnych.
